\section{Main Concepts}

%\subsection{Generalized maps}

\emph{Generalized maps} (or G-maps) \cite{lienhardt1991topological,
  Damiand-Lienhardt14} allow the representation of manifold geometric objects
(with or without boundaries), based on a cellular $n$-dimensional topological
structure. %
The representation of an object as a G-map comes intuitively from its
decomposition into topological cells (vertices, edges, faces, volumes, and so
on). The lowest topological entity in a G-map is called a \emph{dart}. A dart is represented as a subgraph called an \emph{orbit}. 
For example, a volume orbit is noted $\orbit{0,1,2}$, a volume face $\orbit{0,1}$,
a volume edge $\orbit{0,2}$ and a volume vertex $\orbit{1,2}$.

%\subsection{Graph transformation rules}

%Jerboas's \cite{belhaouari2014jerboa} graph transformation rules allow the
%formalization of operations over G-maps. %

Jerboa's \cite{belhaouari2014jerboa} rules describe the transformation of a G-map $G$ into another G-map $H$
by describing a subgraph to match in $G$ and the subgraph replacing it in $H$.
These subgraphs respectively represent the left and right members of a rule. %
A rule's member is made up of nodes describing a G-map's orbit (see Fig.\ref{fig:rule-triangulation}).
%
\begin{figure}[htb]
    \centering
    \includegraphics[height=45pt]{rule-face-triangulation-3D-boxed.pdf}
    \caption{Face triangulation rule}
    \label{fig:rule-triangulation}
\end{figure}
%
%\david{ce serait bien d'insérer la fig de la règle de triangulation avec la petite légende qui va bien pour au moins qu'on voit la tronche d'une règle. Je pense que si la résumé étendu dépasse de quelques lignes ce n'est pas dramatique tant que le tout rentre sur 3 pages.}
Perfor\-ming a syntactical analysis of rules allows to characterize the topological changes
affecting an orbit throughout an application. %
Thus, it is possible to automatically detect topological changes of an orbit and to track it automatically throughout a construction. 
For example, in Fig.\ref{fig:rule-triangulation}, the rule splits the face orbit $\orbit{0,1,3}$ in the green box,
creates the vertex orbit $\orbit{1,2,3}$ in the red box and modifies the volume vertex orbit $\orbit{1,2}$ in the blue box.

%Rules contain the necessary information characterizing the topological changes
%affecting an orbit throughout an application. %
%Thus, the tracking of an orbit is automatically made without any addition other
%than the rules' syntactic analysis. %

\section{Persistent Name}

%During an object's construction, a \emph{parametric specification} records both
%the rules representing the operations applied and their parameters (both
%topological and geometric) in order to describe the modeling process. %
%Since a parametric specification is expected to be edited for future reevaluation,
%topological parameters must be persistent references called persistent names. \david{Cela est déjà dit dans l'intro. Tu peux gagner de la place.}
%Editing a parametric specification means that rules may be added, deleted,
%moved, and their parameters can be changed. %
%With such changes, the topological parameters may have different concrete names,
%be deleted and so on. %
%As a result, using the parameters' concrete names eventually lead to unexpected
%results (at best) or failure at runtime. %


%\subsection{Persistent name}

Since Jerboa's graph transformation rules use darts as topological parameters,
it follows that each persistent name must represent a unique dart. As it
happens, rules make it possible to determine unambiguously which node filters or
creates any dart. %Therefore, a persistent name, noted $PN$, is a list of all the right nodes involved in a darts' history. These nodes are coupled with a rule application's number and sorted.
Therefore, a persistent name, noted $PN$, is a list of all the right nodes involved in a darts' history, coupled with a rule application's number.
%For example, in the $PN_3$ used in the application , the first application's rule node $n_3$ creates a dart $d$, the second application matches $d$ and its rule node $n_4$ creates another dart $d'$. Finally, the third application's matches $d'$.
For example, $PN_3=[1n_{3};2n_{4};3n_{0}]$ in Fig.\ref{fig:fr-color-face}. Indeed, the dart 35 used to designate the face to be colored was created through the nodes $n_3$ of 1-create rule, $n_4$ of 2-extrude rule, and $n_0$ of 3-triangulation rule.
% Since rules use darts as topological parameters, it follows that each persistent
% name must represent a unique dart. %
% As it happens, rules make it possible to determine unambiguously, which node
% filters or creates any dart. %


% \subsection{Parametric Specification syntax}

% A number of fields are used to describe an application within a parametric
% specification, namely an application number, a rule application with
% the topological parameters and the geometric parameters. %

\section{Evaluation}

\begin{figure}[htb]
    \centering
    \begin{subfigure}{.19\columnwidth}
        \includegraphics[width=\textwidth]{figures/fil-rouge-short-create-face.pdf}
        \caption{}
        \label{fig:fr-create-face}
    \end{subfigure}
    \begin{subfigure}{.19\columnwidth}
        \includegraphics[width=\textwidth]{figures/fil-rouge-short-extrude-face.pdf}
        \caption{}
        \label{fig:fr-extrude-face}
    \end{subfigure}
    \begin{subfigure}{.19\columnwidth}
        \includegraphics[width=\textwidth]{figures/fil-rouge-short-triangulate-face.pdf}
        \caption{}
        \label{fig:fr-triangulate-face}
    \end{subfigure}
    \begin{subfigure}{.19\columnwidth}
        \includegraphics[width=\textwidth]{figures/fil-rouge-short-color-face.pdf}
        \caption{}
        \label{fig:fr-color-face}
    \end{subfigure}
    \caption{Evaluation. (\subref{fig:fr-create-face}) 1-create(pos); (\subref{fig:fr-extrude-face}) 2-extrude($PN_1$, vec); (\subref{fig:fr-triangulate-face}) 3-triangulate($PN_2$); (\subref{fig:fr-color-face}) 4-color($PN_3$)}
\end{figure}
%
\begin{figure}[htb]
    \centering
    \begin{subfigure}{.19\columnwidth}
        \includegraphics[width=\textwidth]{figures/fil-rouge-short-create-face.pdf}
        \caption{}
        \label{fig:fr-reeval-create-face}
    \end{subfigure}
     \begin{subfigure}{.19\columnwidth}
        \includegraphics[width=\textwidth]{figures/fil-rouge-short-reeval-insert-vertex.pdf}
        \caption{}
        \label{fig:fr-reeval-insert-vertex}
    \end{subfigure}
    \begin{subfigure}{.19\columnwidth}
        \includegraphics[width=\textwidth]{figures/fil-rouge-short-reeval-extrude-face.pdf}
        \caption{}
        \label{fig:fr-reeval-extrude-face}
    \end{subfigure}
    \begin{subfigure}{.19\columnwidth}
        \includegraphics[width=\textwidth]{figures/fil-rouge-short-reeval-triangulate-face.pdf}
        \caption{}
        \label{fig:fr-reeval-triangulate-face}
    \end{subfigure}
    \begin{subfigure}{.19\columnwidth}
        \includegraphics[width=\textwidth]{figures/fil-rouge-short-reeval-color-face.pdf}
        \caption{}
        \label{fig:fr-reeval-color-face}
    \end{subfigure}
    \caption{Reevaluation. (\subref{fig:fr-create-face}) 1-create(pos); (\subref{fig:fr-reeval-insert-vertex}) ADD1-insert(3); (\subref{fig:fr-reeval-extrude-face}) extrude($PN_1$, vec); (\subref{fig:fr-reeval-triangulate-face}) triangulate($PN_2$); (\subref{fig:fr-reeval-color-face}) 4-color($PN_3$)}
\end{figure}
%Once the initial evaluation of an object is done, a
The evaluation is based on a Directed Acyclic Graph (or DAG) to compute each persistent name before any reevaluation can take place. %
The DAG is built bottom-up by a backward traversal through the persistent name. The role of an evaluation DAG is to trace the history of a topological parameter
back to the first created orbits it originates from, thus allowing matching a corresponding topological parameter at reevaluation. %
\begin{figure}[htb]
  \centering
  \includegraphics[width=0.3\textwidth]{pn7-short}
  \caption{$PN_{3}$'s evaluation DAG}
  \label{fig:eval-dag-pn7}
\end{figure}
In order to accurately represent the history of an orbit, two types of arrows
are used in an evaluation DAG (and later in the reevaluation DAG). A black arrow , or \emph{trace}, 
represents the evolution of an orbit throughout the modeling
process. A red arrow, or \emph{origin}, represents another orbit the tracked orbit
originates from.
For example,
The DAG in Fig.~\ref{fig:eval-dag-pn7} represents the history the topological parameter designated as $PN_3$ in Fig.\ref{fig:fr-color-face} which is the face being colored. This history contains the trace and origin orbits as well as their changes (creation, split, modification, and so on).
%\david{Nous trouvons  dans cet historique, les orbites traces et origines ainsi que l'évolution de ces orbites (création, scission, modification, etc.) } 
%In this history, a square face is created, then is extruded into a cube and, finally, one of the volume's face is triangulated.%

%In Fig.~\ref{fig:eval-dag-pn7}, the $\orbit{0,1}$ at
%orbit level 2 is related to two orbits at orbit level 2, the $\orbit{0,1}$ orbit
%with a trace arrow and the $\orbit{0}$ with an origin arrow. The trace arrow
%relation here means the orbit $\orbit{0,1}$ at orbit level 3 is an evolution of
%the orbit $\orbit{0,1}$ at orbit level 2. The origin arrow is added to reflected
%that since the orbit $\orbit{0,1}$ is split by the \texttt{3-triangulation} rule 
%application, the orbit tracked in the DAG comes from the orbit $\orbit{0}$ at orbit level 2.


\section{Reevaluation}

%When reevaluating, editing
Since topological parameters may change after editing a parametric specification, % makes the topological parameters subject to changes. %
it is necessary to build reevaluation DAGs from evaluation DAGs in
order to update topological parameters. %
% Once built,
Reevaluation DAGs can designate one, several, or no orbit depending
on the editing of the parametric specification.

%\subsection{Reevaluation DAG}

%Contrary to the related evaluation DAG, a
Reevaluation DAGs are built top-down. While an evaluation DAG represents the orbit’s history of a topological parameter, the
reevaluation DAG represents the history of this same orbit after
editing the parametric specification.
\begin{figure}[htb]
    \centering
    \includegraphics[width=.95\columnwidth]{pn7-reval-short}
    \caption{$PN_3$'s reevaluation DAG}
    \label{fig:pn7-reval-dag}
\end{figure}
%\subsection{Parameter matching strategies}
Such an edit yields a different DAG at
reevaluation (with event levels and/or branches being added, deleted or both).
% [Xavier] la phrase ci-dessous semble redondante par rapport à la précédente, autant l'enlever pour gagner de la place
%For example, an orbit split present in the evaluation DAG may disappear during the reevaluation, a merging can be added and so on.
Several matching strategies
can then be considered depending on the application’s context, such as the DAG in Fig.~\ref{fig:pn7-reval-dag}
representing the histories of the faces to color in blue in Fig.\ref{fig:fr-reeval-color-face}.

\section{Conclusion}
This paper introduces our work on extending the scope of naming problem studies to encompass rule-based graph transformation modeling systems.
We implement a persistent naming scheme identifying a unique dart's history. % through the successive operations and their specific node being involved.
We represent any topological parameter's history with an evaluation DAG from a persistent name.
We can match one, several or no topological parameter depending on the editing to reevaluate.

More complex operations can make use of several rules brought
together in a script. Later works will revolve around widening this
reevaluation mechanism to scripts.

%%% Local Variables:
%%% mode: latex
%%% TeX-master: "../main"
%%% End:
